\chapter{GHB (Gamma-Hydroxybutyrate)}
\index{GHB (Gamma-Hydroxybutyrate)}

\section*{Definition and Overview}
GHB (gamma-hydroxybutyrate) is a central nervous system depressant drug that is used illicitly in many countries, often as a liquid or powder. At low doses it produces relaxation, euphoria, and increased sociability; at higher doses it causes sedation, confusion, and loss of consciousness. GHB has a narrow margin between the dose that produces the desired effect and the dose that causes overdose, and its effects come on rapidly, which makes it one of the most dangerous drugs used recreationally. It is also associated with drug-facilitated sexual assault because it is odourless, colourless, and tasteless and causes amnesia. GHB is sometimes used medically in controlled settings overseas for narcolepsy, but it has no legal therapeutic use in many countries.

\section*{Epidemiology}
GHB is used by a relatively small but significant proportion of people who use drugs in many countries, particularly within party and nightlife scenes and among some men who have sex with men. Use is concentrated in major cities and among people in their 20s and 30s. Because GHB is not detected in routine drug tests and overdoses often present as collapses at venues, the true scale of use and harm is likely underestimated. GHB is involved in a disproportionate number of drug-related emergency department presentations and deaths relative to the number of people who use it, reflecting its narrow safety margin and rapid onset.

\section*{Aetiology and Pathophysiology}
GHB acts on the brain by binding to GABA receptors, the main inhibitory neurotransmitter system.

\begin{itemize}
    \item \textbf{Mechanism:} GHB enhances GABA activity, slowing brain function and producing sedation, relaxation, and, at higher doses, respiratory depression and unconsciousness
    \item \textbf{Rapid absorption:} effects begin within 15--30 minutes of ingestion and peak quickly
    \item \textbf{Short duration:} the acute effects usually last 3--6 hours, but the window between intoxication and overdose is narrow
    \item \textbf{Interactions:} alcohol and other depressants markedly increase the risk of overdose and death
    \item \textbf{Dependence:} regular use leads to tolerance and physical dependence; withdrawal can be severe and life-threatening
    \item \textbf{Amnesia:} at sedating doses, GHB impairs memory formation, which is why it is used in drug-facilitated assault
\end{itemize}

\section*{Clinical Features}
The effects of GHB depend heavily on the dose and on whether other drugs are present.

\begin{itemize}
    \item \textbf{Low doses:} relaxation, euphoria, reduced inhibitions, increased talkativeness, and mild dizziness
    \item \textbf{Moderate doses:} drowsiness, slurred speech, unsteady gait, nausea, and confusion
    \item \textbf{High doses:} loss of consciousness, slowed breathing, bradycardia, vomiting, and coma
    \item \textbf{Overdose:} respiratory depression, seizure, and cardiac arrest can be fatal
    \item \textbf{On withdrawal:} anxiety, insomnia, tremor, sweating, agitation, and in severe cases seizures, hallucinations, and delirium
    \item \textbf{Amnesia:} gaps in memory for events while intoxicated
\end{itemize}

\section*{Differential Diagnosis}
A person presenting unconscious or acutely confused after suspected GHB use should be assessed for other causes:

\begin{itemize}
    \item Alcohol and other depressant drug overdose, including opioids and benzodiazepines
    \item Other drugs of misuse, including methamphetamine, cocaine, or novel psychoactive substances
    \item Head injury or intracranial bleeding
    \item Hypoglycaemia, seizures, and metabolic disorders
    \item Meningitis or encephalitis
    \item Stroke and other acute neurological events
    \item Drug-facilitated sexual assault, which should be considered and addressed with appropriate support
\end{itemize}

\section*{When to Seek Medical Care}
Any collapse, unconsciousness, or severe drowsiness after drug use is a medical emergency requiring immediate attention. People who use GHB regularly should discuss their use with a health professional, as dependence and withdrawal require medical management.

\textbf{Contact your local emergency services for:}
\begin{itemize}
    \item Loss of consciousness or an inability to wake someone
    \item Slow, shallow, or irregular breathing
    \item Seizures or convulsions
    \item Vomiting in an unconscious person
    \item Agitation, confusion, or signs of withdrawal in a regular user
    \item Any suspected drug-facilitated assault
\end{itemize}

\section*{Diagnosis and Investigations}
Diagnosis is largely clinical and supportive, guided by the circumstances and the person's presentation.

\begin{itemize}
    \item \textbf{History and witness account:} details of what was taken, how much, and when
    \item \textbf{Clinical assessment:} level of consciousness, breathing, heart rate, and pupil size
    \item \textbf{Blood glucose and pulse oximetry:} immediate bedside checks to exclude other causes
    \item \textbf{Toxicology:} GHB is cleared from the body within hours and is not detected by standard drug tests; specific laboratory testing may be arranged in forensic settings
    \item \textbf{Monitoring:} people with significant overdose are observed in hospital until they recover
    \item \textbf{Assessment for assault:} where drug-facilitated assault is suspected, specialist services provide forensic examination and support
\end{itemize}

\section*{Management}
Management focuses on emergency stabilisation, supportive care, and treatment of dependence.

\begin{itemize}
    \item \textbf{Emergency care:} airway support, breathing assistance, and monitoring in hospital; GHB overdose is managed supportively because there is no antidote
    \item \textbf{Airway protection:} unconscious people are placed on their side and may need intubation if breathing is compromised
    \item \textbf{Withdrawal management:} medical detoxification under supervision, often in hospital, because GHB withdrawal can be dangerous
    \item \textbf{Support for dependence:} counselling, relapse prevention, and referral to drug and alcohol services
    \item \textbf{Harm reduction:} honest information about the risks of GHB and strategies to reduce harm
    \item \textbf{Support after assault:} counselling, crisis services, and police options for people affected by drug-facilitated assault
\end{itemize}

\section*{Complications}
\begin{itemize}
    \item Fatal overdose from respiratory depression
    \item Aspiration of vomit while unconscious
    \item Seizures, coma, and brain injury from hypoxia
    \item Serious physical injuries from falls and accidents while intoxicated
    \item Drug-facilitated sexual assault and its psychological consequences
    \item Dependence with severe, potentially life-threatening withdrawal
    \item Long-term effects on memory, mood, and mental health
    \item Legal consequences of possession and supply
\end{itemize}

\section*{Prognosis and Outcomes}
With prompt emergency care, most people survive individual GHB overdoses, and those who are treated in hospital usually recover within hours. The larger dangers are the cumulative risks of regular use: dependence, withdrawal, overdose, and trauma. People who seek treatment for GHB dependence can achieve stable recovery with supervised withdrawal and ongoing psychosocial support. For anyone affected by GHB-related assault, specialised support services can help with recovery and legal options.

\section*{Prevention and Self-Care}
\begin{itemize}
    \item Understand that GHB has an extremely narrow margin between the desired effect and overdose
    \item Never combine GHB with alcohol or other depressant drugs, which dramatically increases the risk of death
    \item Avoid using alone, so someone can call for help if you collapse
    \item Never accept drinks you have not seen being poured, and keep your drink covered, to reduce the risk of drink spiking
    \item Learn the signs of overdose and how to seek emergency help
    \item If you use drugs, seek honest advice from services such as Alcohol and Drug Information Service or harm reduction organisations
    \item If you think you are dependent, seek medical help for safe withdrawal rather than stopping abruptly
    \item If you have been assaulted while under the influence of a drug, contact services such as 1800RESPECT for support
\end{itemize}

\section*{Sources and Further Reading}
The following reputable sources provide further information for healthcare students and patients seeking reliable, current advice about GHB.

\begin{itemize}
    \item Alcohol and Drug Foundation. \textit{GHB: effects, risks, and harms.} https://adf.org.au/
    \item Healthdirect Australia. \textit{GHB (gamma-hydroxybutyrate).} https://www.healthdirect.gov.au/
    \item Australian Government Department of Health and Aged Care. \textit{Illicit drugs: GHB information.} https://www.health.gov.au/
    \item Penington Institute. \textit{GHB: safety and overdose information.} https://www.penington.org.au/
    \item 1800RESPECT. \textit{Support after sexual assault.} https://www.1800respect.org.au/
    \item NHS. \textit{GHB: the risks of gamma-hydroxybutyrate.} https://www.nhs.uk/
\end{itemize}
